% GENERATED FILE. Edit canonical lessons, then run npm run generate:books.
% canonical-source-hash: fnv1a64:82350d560f45f362
% canonical-lessons: TE-C132-saksu, TE-C132-kotu, TE-C132-dhoti, TE-C132-jakettu, TE-C132-tamarapuvvu

\chapter{Socks, Coat, Dhoti, Blouse, Lotus}
\label{ch:te-socks-coat-dhoti-blouse-lotus}
\hlchaptermodality{132}

\begin{chapteropening}
\textbf{By the end of this chapter:} \emph{I can say socks, a coat, a dhoti, a blouse and a lotus.}

Complete the last lesson of chapter 132: i can say socks, a coat, a dhoti, a blouse and a lotus.
\end{chapteropening}

\section[\texorpdfstring{\te{సాక్సు} (sāksu)}{సాక్సు (sāksu)}]{\te{సాక్సు} (sāksu) — socks}

\label{lesson:TE-C132-saksu}



\begin{quote}
Before the new one: say the Telugu for a pocket, then the Telugu for a belt.
\end{quote}

\subsection*{You'll want to know: \te{సాక్సు}}
\textbf{\te{సాక్సు}} — \emph{sāksu} — \textquotedblleft{}socks\textquotedblright{}.

\begin{grammarlens}[title={What you've built}]
One more word for this chapter.
\end{grammarlens}

\subsection*{Your turn}
\begin{itemize}
\raggedright
  \item \emph{Say it:} \emph{sāksu}
  \item \emph{Say it:} \emph{sāksu}, once more
  \item \emph{Say it:} say \emph{sāksu}
  \item \emph{Recall:} say the Telugu for a pocket, then the Telugu for a belt, then say \emph{sāksu} again
\end{itemize}

\begin{tcolorbox}[breakable,colback=teal!4,colframe=teal!35!black,title={Before you move on}]
What does \te{సాక్సు} mean? (\textquotedblleft{}Socks\textquotedblright{}.) Say it once more.
\end{tcolorbox}
\section[\texorpdfstring{\te{కోటు} (kōṭu)}{కోటు (kōṭu)}]{\te{కోటు} (kōṭu) — a coat}

\label{lesson:TE-C132-kotu}



\begin{quote}
Before the new one: say the Telugu for a belt, then the Telugu for socks.
\end{quote}

\subsection*{You'll want to know: \te{కోటు}}
\textbf{\te{కోటు}} — \emph{kōṭu} — \textquotedblleft{}a coat\textquotedblright{}.

\begin{grammarlens}[title={What you've built}]
One more word for this chapter.
\end{grammarlens}

\subsection*{Your turn}
\begin{itemize}
\raggedright
  \item \emph{Say it:} \emph{kōṭu}
  \item \emph{Say it:} \emph{kōṭu}, once more
  \item \emph{Say it:} say \emph{kōṭu}
  \item \emph{Recall:} say the Telugu for a belt, then the Telugu for socks, then say \emph{kōṭu} again
\end{itemize}

\begin{tcolorbox}[breakable,colback=teal!4,colframe=teal!35!black,title={Before you move on}]
What does \te{కోటు} mean? (\textquotedblleft{}A coat\textquotedblright{}.) Say it once more.
\end{tcolorbox}
\section[\texorpdfstring{\te{ధోతి} (dhōti)}{ధోతి (dhōti)}]{\te{ధోతి} (dhōti) — a dhoti}

\label{lesson:TE-C132-dhoti}



\begin{quote}
Before the new one: say the Telugu for socks, then the Telugu for a coat.
\end{quote}

\subsection*{You'll want to know: \te{ధోతి}}
\textbf{\te{ధోతి}} — \emph{dhōti} — \textquotedblleft{}a dhoti\textquotedblright{}.

\begin{grammarlens}[title={What you've built}]
One more word for this chapter.
\end{grammarlens}

\subsection*{Your turn}
\begin{itemize}
\raggedright
  \item \emph{Say it:} \emph{dhōti}
  \item \emph{Say it:} \emph{dhōti}, once more
  \item \emph{Say it:} say \emph{dhōti}
  \item \emph{Recall:} say the Telugu for socks, then the Telugu for a coat, then say \emph{dhōti} again
\end{itemize}

\begin{tcolorbox}[breakable,colback=teal!4,colframe=teal!35!black,title={Before you move on}]
What does \te{ధోతి} mean? (\textquotedblleft{}A dhoti\textquotedblright{}.) Say it once more.
\end{tcolorbox}
\section[\texorpdfstring{\te{జాకెట్టు} (jākeṭṭu)}{జాకెట్టు (jākeṭṭu)}]{\te{జాకెట్టు} (jākeṭṭu) — a blouse}

\label{lesson:TE-C132-jakettu}



\begin{quote}
Before the new one: say the Telugu for a coat, then the Telugu for a dhoti.
\end{quote}

\subsection*{You'll want to know: \te{జాకెట్టు}}
\textbf{\te{జాకెట్టు}} — \emph{jākeṭṭu} — \textquotedblleft{}a blouse\textquotedblright{}.

\begin{grammarlens}[title={What you've built}]
One more word for this chapter.
\end{grammarlens}

\subsection*{Your turn}
\begin{itemize}
\raggedright
  \item \emph{Say it:} \emph{jākeṭṭu}
  \item \emph{Say it:} \emph{jākeṭṭu}, once more
  \item \emph{Say it:} say \emph{jākeṭṭu}
  \item \emph{Recall:} say the Telugu for a coat, then the Telugu for a dhoti, then say \emph{jākeṭṭu} again
\end{itemize}

\begin{tcolorbox}[breakable,colback=teal!4,colframe=teal!35!black,title={Before you move on}]
What does \te{జాకెట్టు} mean? (\textquotedblleft{}A blouse\textquotedblright{}.) Say it once more.
\end{tcolorbox}
\section[\texorpdfstring{\te{తామరపువ్వు} (tāmarapuvvu)}{తామరపువ్వు (tāmarapuvvu)}]{\te{తామరపువ్వు} (tāmarapuvvu) — a lotus}

\label{lesson:TE-C132-tamarapuvvu}



\begin{quote}
Before the new one: say the Telugu for a dhoti, then the Telugu for a blouse.
\end{quote}

\subsection*{You'll want to know: \te{తామరపువ్వు}}
\textbf{\te{తామరపువ్వు}} — \emph{tāmarapuvvu} — \textquotedblleft{}a lotus\textquotedblright{}.

\begin{grammarlens}[title={What you've built}]
One more word for this chapter.
\end{grammarlens}

\subsection*{Your turn}
\begin{itemize}
\raggedright
  \item \emph{Say it:} \emph{tāmarapuvvu}
  \item \emph{Say it:} \emph{tāmarapuvvu}, once more
  \item \emph{Say it:} say \emph{tāmarapuvvu}
  \item \emph{Recall:} say the Telugu for a dhoti, then the Telugu for a blouse, then say \emph{tāmarapuvvu} again
\end{itemize}

\begin{tcolorbox}[breakable,colback=teal!4,colframe=teal!35!black,title={Before you move on}]
What does \te{తామరపువ్వు} mean? (\textquotedblleft{}A lotus\textquotedblright{}.) Say it once more.
\end{tcolorbox}
